\subsection{Removing all single-prime word images from the quadratic question}

Here a pro-odd group means a profinite group all of whose finite
continuous quotients have odd order. This term does not require
commutativity.

\begin{lemma}\label{lem:relative-proodd-restriction}
Let \(G=QH\), where \(Q\) is a closed normal subgroup of the
profinite group \(G\) and \(H\) is closed. Let \(E\subset Q\)
be closed and normal in \(G\), and suppose that \(Q/E\) is pro-odd.
For every finite abelian \(2\)-group \(A\), restriction induces
a bijection
\[
 \begin{aligned}
 \{\chi\in\operatorname{Hom}_{\rm cts}(G,A):\chi|_H=0\}
 &\longrightarrow\\
 &\{\psi\in\operatorname{Hom}_{\rm cts}(E,A)^G:
                  \psi|_{E\cap H}=0\}.
 \end{aligned}
\]
All maps on the left are continuous homomorphisms. In particular
there is a canonical isomorphism
\[
 \left(Q/\overline{[G,Q](Q\cap H)}\right)_{\max2}
 \ \simeq\
 \left(E/\overline{[G,E](E\cap H)}\right)_{\max2}.
\]
No containment \(Q\cap H\subset E\) is required.
\end{lemma}
\begin{proof}
Given \(\psi\) on the right, put \(L=\ker\psi\).
Then \(L\) is closed and normal in \(G\), and \(E/L\) is a
finite central \(2\)-subgroup of \(Q/L\), whose quotient is pro-odd.
We first record why a central finite \(2\)-subgroup \(C\) of a
finite group \(J\), with \(J/C\) of odd order, has a unique normal
complement. The transfer \(V:J\to C\) is a homomorphism and
satisfies \(V(c)=c^{[J:C]}\) for \(c\in C\). These assertions
follow directly by multiplying the coset-representative factors;
centrality permits their reordering. Raising \(V\) to an exponent
inverse to \([J:C]\) modulo the exponent of \(C\) gives a
retraction \(J\to C\). Its kernel is an odd normal complement.
Thus \(J=C\times P\). Any other complement is the graph of a
homomorphism from the odd group \(P\) to \(C\), which must be
zero. The complement is unique.

Apply this observation to every finite quotient of \(Q/L\).
The unique odd complements are compatible with all quotient maps.
Taking inverse limits gives the canonical direct product
\[
                    Q/L=(E/L)\times P
\]
with \(P\) pro-odd. Projection to \(E/L\), followed by \(\psi\),
extends \(\psi\) to a continuous homomorphism \(\widetilde\psi:Q\to A\).
Any two extensions differ by a homomorphism \(Q/E\to A\),
which is zero. Uniqueness implies that \(\widetilde\psi\) is
\(G\)-invariant.

The restriction of \(\widetilde\psi\) to \(Q\cap H\) vanishes on
\(E\cap H\). The quotient
\[
 (Q\cap H)/(E\cap H)\ \hookrightarrow\ Q/E
\]
is pro-odd, so this restriction is zero. Hence
\(\chi(qh)=\widetilde\psi(q)\) is well-defined.
Invariance proves that it is a homomorphism, and the compact
surjection \(Q\times H\to G\) proves continuity.
The same pro-odd argument gives uniqueness and injectivity.
Applying the bijection to all finite abelian \(2\)-groups proves
the asserted canonical isomorphism of maximal pro-two quotients.
\end{proof}

Return to \(G=\widehat{\mathrm{GT}}\), \(H=G_{\mathbf Q}\),
\(D=N_{\rm rad}\), and the established normalization
\(G=R_DH\), where \(R_D\) is the kernel on \(D\) of the
free pro-two word action. For every prime \(p\), let
\[
 \Phi_p:G\longrightarrow
       \operatorname{Aut}(\widehat F_2^{(p)})
\]
be the standard action on the free pro-\(p\) quotient, and put
\[
 P_D=D\cap\bigcap_p\ker\Phi_p,\qquad
 T_D=P_D\cap\ker d_{(3)}.
\]
On multiplier one, \(\Phi_p(g)=1\) is equivalent to \(f_g=1\)
in the free pro-\(p\) group. Indeed the normalized action fixes
\(x\), and triviality on \(y\) says that \(f_g\) centralizes
the free generator \(y\). Its procyclic centralizer intersects
the commutator subgroup trivially.

\begin{theorem}\label{thm:all-prime-residual-quadratic-domain}
The subgroups \(P_D,T_D\) are closed and normal in \(G\).
Both \(R_D/P_D\) and \(R_D/T_D\) are pro-odd. Consequently
\[
 \begin{aligned}
 B_{\max2}&\simeq
 \left(P_D/\overline{[G,P_D](P_D\cap H)}\right)_{\max2}\\
 &\simeq
 \left(T_D/\overline{[G,T_D](T_D\cap H)}\right)_{\max2}.
 \end{aligned}
\]
Restriction from \(X_2(N)\) to either \(P_D\) or \(T_D\)
has the same kernel as restriction to \(R_D\), namely
\[
 \left(\bigoplus_{q\ {\rm odd\ prime}}\mathbf F_2\psi_q\right)
                         \oplus\mathbf F_2\eta_3.
\]
These assertions concern the maximal pro-two relative quotient;
they do not assert \(G=P_DH\) or \(G=T_DH\).
\end{theorem}
\begin{proof}
For a finitely generated free pro-\(p\) group, the kernel of
its automorphism action on its Frattini quotient is pro-\(p\).
For completeness, use the lower exponent-\(p\) central series
\(V_1=F\), \(V_{n+1}=V_n^p[V_n,F]\).
The \(V_n\) are characteristic open subgroups forming a
neighborhood basis. At each successive finite quotient,
automorphisms that are the identity one step earlier are
determined by central corrections to the free generators in
\(V_n/V_{n+1}\); their group is a subgroup of a finite
elementary abelian \(p\)-group. Induction, beginning with the
Frattini quotient, and inverse limits prove the assertion.

Every element of \(D\) has multiplier one. Its pro-\(p\)
action is therefore in this pro-\(p\) automorphism kernel.
Since the pro-two action is already trivial on \(R_D\),
the quotient \(R_D/P_D\) embeds as a closed subgroup in a
product of pro-\(p\) groups over odd primes \(p\).
Every finite quotient of that product, and of its closed
subgroups, has odd order. Thus \(R_D/P_D\) is pro-odd.
The quotient \(R_D/(R_D\cap\ker d_{(3)})\) has exponent three.
The diagonal map embeds \(R_D/T_D\) in the product of these
two pro-odd quotients, proving the second assertion.
Normality follows from the characteristic nature of the
free pro-\(p\) quotients and the equivariance of \(d_{(3)}\).

Apply Lemma~\ref{lem:relative-proodd-restriction} with \(Q=R_D\).
For the original character kernel, a character on \(N\) that
vanishes on \(P_D\), or on \(T_D\), restricts to a homomorphism
from the corresponding pro-odd quotient of \(R_D\) to
\(\mathbf F_2\). It is therefore zero on \(R_D\).
The reverse implication is immediate. The known kernel for
restriction to \(R_D\) completes the proof.
\end{proof}

\begin{proposition}\label{prop:all-prime-residual-strictness}
Both \(P_D\) and \(T_D\) have infinite index in \(R_D\).
Moreover \(T_D\) has infinite index in \(R_D\cap\ker d_{(3)}\).
\end{proposition}
\begin{proof}
Fix any odd prime \(p\). We use only the existence of a
nonzero arithmetic degree-three class, not arithmetic
saturation at \(p\): Ihara's nonzero Soul\'e class is recalled
in \cite[Remark~8.6, p.~151]{HainMatsumoto2003}.
Choose an arithmetic element \(\sigma\) with \(\lambda_p(\sigma)=1\)
and nonzero degree-three pro-\(p\) leading term \(q\).
Here the passage from the outer-action statement to the
normalized word has no kernel in this degree. The normalized
word has zero degrees one and two when \(\lambda_p=1\).
Its degree-three term \(q\) induces the derivation
\(S\mapsto0,\ T\mapsto[T,q]\), up to the harmless convention
of reversing this bracket. An inner derivation with this
property would be given by a degree-three Lie element
commuting with \(S\); the centralizer of \(S\) in the free
Lie algebra has no degree-three part. Thus a nonzero outer
degree-three class has \(q\ne0\). Multiplying an integral
representative if necessary realizes the nonzero rational
class by a Galois element.

Choose \(h\in H\) with \(u=\lambda_p(h)=1+p\), and use the
commutator convention \([a,b]=aba^{-1}b^{-1}\). Put
\[
 z=[h,\sigma],\qquad
 d=hzh^{-1}z^{-\lambda(h)}.
\]
The exponent in the second expression is the full adelic
integer \(-\lambda(h)\): the continuous power map
\(\widehat{\mathbf Z}\to\overline{\langle z\rangle}\)
defines it. The multiplier of \(z\) is one at every prime.
The degree-three word coordinate is additive on multiplier
one, and conjugation has weight \(\lambda_p^3\). Hence
\[
 \operatorname{in}_3(z)=(u^3-1)q,\qquad
 \operatorname{in}_3(d)=(u^3-u)(u^3-1)q\ne0
                  \quad\text{over }\mathbf Z_p.
\]
Each rational Kummer coordinate is additive on multiplier
one and has conjugation weight \(\lambda(h)\).
It therefore vanishes on \(d\), so \(d\in D\cap H\).

Let \(\epsilon_p\in\widehat{\mathbf Z}=\prod_r\mathbf Z_r\)
have component one at \(p\) and zero at every other prime.
The element \(w=d^{\epsilon_p}\) is still in \(D\cap H\).
The pro-two image of \(d\) lies in a pro-two automorphism
group, so the zero two-component of \(\epsilon_p\) kills
that image. Thus \(w\in R_D\cap H\).
Its degree-three pro-\(p\) term equals that of \(d\), and is
nonzero. Since this term is an additive coordinate in a
torsion-free \(\mathbf Z_p\)-module, the pro-\(p\) image
of the closed cyclic subgroup generated by \(w\) is infinite.
The subgroup \(P_D\) kills this image, proving infinite index.
Finally \(d_{(3)}\) is zero on \(H\), so \(w\) already lies in
\(R_D\cap\ker d_{(3)}\); no additional cubing is needed.
The same infinite pro-\(p\) image proves the assertion for
\(T_D\).
\end{proof}
